% Evaluación de Filosofía — Décimo — Cuarto periodo 2026
% Tema: de la Antigüedad al Renacimiento. Herencia de Platón y Aristóteles, san Agustín, Tomás de
% Aquino y la escolástica, la «Edad Oscura», Renacimiento y humanismo (Erasmo), la Utopía de Moro.
% Compilar: bash .claude/skills/evaluaciones-latex/scripts/compilar.sh evaluaciones/ejemplos/2026-11-14_filosofia-10.tex
\documentclass[carta,doscolumnas,compacta]{../plantilla/evaluacion}
\usepackage{../plantilla/evaluacion-extras}

% Colegio, logo, color y docente: de curriculo/config/colegio.yaml y docente.yaml.
% Para otro valor en este documento: \colegio{…} \logo{archivo.png} \colorcolegio{166534} \docente{…}

\asignatura{Filosofía}
\grado{Décimo}
\periodo{Cuarto periodo}
\anio{2026}
\fecha{14/11/2026}
\titulo{Evaluación de Filosofía: de la escolástica a la Utopía}
\instrucciones{Tiempo: 1 hora. Todo lo que necesitas para responder está en los textos, fragmentos,
esquemas e imágenes: léelos con atención antes de contestar. En las preguntas de selección marca
una sola respuesta (A, B, C o D). En las de \textbf{afirmación y razón} decide si cada parte es
verdadera y si la razón explica la afirmación. En las abiertas se califica la claridad y la solidez
de tus argumentos.}

\begin{document}
\encabezado

\begin{preguntas}

% ================================================================ Bloque 1: herencia griega
\begin{contexto}[Una larga herencia]
\grupo{Responde las preguntas 1 a 3 con base en el texto, la imagen y la línea de tiempo.}
La filosofía medieval no partió de cero. San Agustín leyó a Platón a través de sus seguidores, los
neoplatónicos. Las obras de Aristóteles, en cambio, se habían perdido en gran parte en Europa
occidental y volvieron en los siglos XII y XIII gracias a traducciones del árabe y del griego al
latín, como las de la escuela de traductores de Toledo y los comentarios del filósofo andalusí
Averroes. Sobre ellas trabajó Tomás de Aquino.
\imagenmarcada[0.98]{img/comun/atenas.jpg}{Rafael, \emph{La escuela de Atenas}
(1509--1511), Palacios Vaticanos. Platón (1) señala hacia arriba; Aristóteles (2) extiende la mano
con la palma hacia la tierra.}{0.475/0.63/1, 0.555/0.63/2}
\lineadetiempo[desde=-450,hasta=1550,paso=250,alto=0.42]
  {-427/Platón/1, -384/Aristóteles/2.2, 354/Agustín/1, 476/{caída de Roma}/-1,
   1088/{1.ª universidad}/-1, 1225/{Tomás de Aquino}/1.8, 1450/imprenta/-1.9, 1516/\emph{Utopía}/1}
\end{contexto}

\begin{pregunta}[D]
Rafael pintó a los filósofos griegos en el palacio del papa, en Roma, hacia 1510. ¿Qué muestra
este hecho sobre el Renacimiento?
\opcioneslargas
  {Que Platón y Aristóteles todavía vivían en el siglo XVI y posaron para el pintor en Roma.}
  {Que la Iglesia prohibía la filosofía griega y Rafael la pintó para desafiar al papa.}
  {Que el Renacimiento rechazaba todo lo cristiano y quería volver a la religión griega.}
  {Que hubo un gran interés por recuperar la cultura clásica griega y romana, incluso en Roma.}
\justificacion{El encargo papal de una obra que honra a los filósofos griegos expresa el interés
  renacentista por la Antigüedad. A es un anacronismo; B contradice que la obra esté en el
  Vaticano; C exagera: el Renacimiento combinó lo clásico y lo cristiano.}
\especificacion{Interpretativa}{Interpreta una obra de arte como documento histórico.}
\end{pregunta}

\begin{pregunta}[A]
Observa los gestos de Platón y Aristóteles en la pintura. ¿Cuál relación con los pensadores
medievales es coherente con el texto y la imagen?
\opcioneslargas
  {Agustín se apoyó en Platón, que mira hacia lo que está más allá de lo sensible; Tomás, en Aristóteles, que parte de este mundo.}
  {Agustín se apoyó en Aristóteles, que mira hacia el cielo; Tomás, en Platón, que mira hacia la tierra y la naturaleza.}
  {Agustín y Tomás rechazaron por completo a los griegos, porque eran filósofos anteriores al cristianismo.}
  {Tomás de Aquino tradujo él mismo todas las obras de Aristóteles del árabe al latín en Toledo.}
\justificacion{Platón (gesto hacia arriba) inspira a Agustín; Aristóteles (gesto hacia la tierra)
  a Tomás. B invierte autores y gestos; C contradice el texto; D le atribuye a Tomás el trabajo de
  los traductores.}
\especificacion{Argumentativa}{Relaciona símbolos visuales con la herencia filosófica.}
\end{pregunta}

\begin{pregunta}[C]
Según la línea de tiempo, lee las afirmaciones:
\begin{afirmaciones}
\item Entre el nacimiento de Aristóteles y el de Tomás de Aquino pasaron más de 1500 años.
\item Agustín vivió en la misma época en que se fundaron las primeras universidades.
\item La \emph{Utopía} de Moro se publicó cuando ya existía la imprenta.
\end{afirmaciones}
¿Cuáles son correctas?
\opciones{Solo I}{Solo II}{I y III}{II y III}
\justificacion{I: de 384 a.\,C. a 1225 d.\,C. hay unos 1600 años. III: 1516 es posterior a 1450.
  II es falsa: Agustín murió en el 430 y la primera universidad es de 1088.}
\especificacion{Interpretativa}{Calcula y compara distancias en una línea de tiempo.}
\end{pregunta}

% ================================================================ Bloque 2: san Agustín
\begin{contexto}[San Agustín de Hipona (354--430)]
\grupo{Responde las preguntas 4 y 5 con base en los textos.}
\retrato[0.27]{img/filosofia10/agustin.jpg}{San Agustín, grabado (s.\ XVIII). Lleva el corazón en
llamas, símbolo de su búsqueda.}%
{\small Agustín nació en el norte de África, en el Imperio romano. Tuvo una juventud agitada,
estudió retórica y filosofía, y se convirtió al cristianismo a los 32 años. Fue obispo de Hipona.
Cuando los visigodos saquearon Roma (410), escribió \emph{La ciudad de Dios}.}
\fragmento{Nos hiciste, Señor, para ti, y nuestro corazón está inquieto hasta que descanse en
ti.}{Agustín, \emph{Confesiones}, I, 1}
\fragmento{No quieras salir fuera de ti; vuelve a ti mismo: en el interior del hombre habita la
verdad.}{Agustín, \emph{La verdadera religión}, 39, 72}
\fragmento{Comprende para creer, cree para comprender.}{Agustín, \emph{Sermón} 43, 9}
\end{contexto}

\begin{pregunta}[B]
Según el segundo fragmento, ¿dónde debe buscar la verdad una persona?
\opcioneslargas
  {En los libros de los filósofos griegos, leyéndolos con mucho cuidado y sin dudar de ellos.}
  {En su propio interior, en su alma, más que en el mundo exterior que perciben los sentidos.}
  {En la observación cuidadosa de la naturaleza, como hacían los primeros filósofos griegos.}
  {En las leyes del Imperio romano, que guardaban la sabiduría de muchos siglos de historia.}
\justificacion{«Vuelve a ti mismo […] en el interior del hombre habita la verdad». A, C y D
  ponen la verdad fuera de la persona, que es justo lo que el fragmento desaconseja.}
\especificacion{Interpretativa}{Identifica la tesis de un fragmento breve.}
\end{pregunta}

\begin{pregunta}[D]
¿Qué relación entre fe y razón expresa el tercer fragmento?
\opcioneslargas
  {Que la razón debe abandonarse, porque basta con creer para llegar a la verdad.}
  {Que solo la razón conduce a la verdad, y que la fe es un obstáculo que conviene superar.}
  {Que la fe y la razón son contrarias, y por eso cada persona debe elegir una de las dos.}
  {Que la fe y la razón se ayudan: la fe orienta la búsqueda y la razón ayuda a entender lo que se cree.}
\justificacion{La frase tiene dos direcciones: comprender para creer y creer para comprender.
  A y B eliminan una de las dos; C las opone, lo contrario de lo que dice el fragmento.}
\especificacion{Interpretativa}{Interpreta una tesis sobre la relación entre fe y razón.}
\end{pregunta}

% ================================================================ Bloque 3: Tomás de Aquino
\begin{contexto}[Tomás de Aquino (1225--1274) y la escolástica]
\grupo{Responde las preguntas 6 a 8 con base en el texto y el esquema.}
\retrato[0.25]{img/filosofia10/tomas.jpg}{Tomás de Aquino, grabado. El sol en el pecho simboliza
la luz de su saber.}%
{\small La \textbf{escolástica} fue la filosofía que se enseñaba en las escuelas y universidades
medievales. Su gran obra es la \emph{Suma teológica} de Tomás de Aquino, un dominico italiano que
enseñó en la Universidad de París. Cada uno de sus cientos de «artículos» sigue el mismo método:}
\begin{center}
\begin{tikzpicture}[font=\fontsize{6.5}{7.8}\selectfont,node distance=2.2mm,
  p/.style={draw=acento,fill=white,rounded corners=1.5pt,text width=0.86\linewidth,inner sep=2.5pt,align=left}]
  \node[p] (a) {\textbf{Pregunta.} \emph{¿Existe Dios?} (\emph{Suma}, I, cuestión 2, artículo 3)};
  \node[p,below=of a] (b) {\textbf{«Parece que no».} Objeciones: las mejores razones en contra (por ejemplo: si Dios existiera, no existiría el mal).};
  \node[p,below=of b] (c) {\textbf{«Pero en contra».} Una autoridad que apoya la otra postura (la Biblia, un filósofo).};
  \node[p,below=of c,fill=acento!12] (d) {\textbf{«Respondo».} La respuesta propia, argumentada (aquí, las «cinco vías»).};
  \node[p,below=of d] (e) {\textbf{Respuestas a las objeciones.} Se responde una por una a las razones del principio.};
  \foreach \x/\y in {a/b,b/c,c/d,d/e} \draw[flecha] (\x) -- (\y);
\end{tikzpicture}
\end{center}
\end{contexto}

\begin{pregunta}[D]
¿Por qué Tomás empieza cada artículo exponiendo las razones \textbf{contrarias} a su propia postura?
\opcioneslargas
  {Porque no estaba seguro de lo que pensaba y prefería que el lector decidiera por sí mismo.}
  {Porque las objeciones eran sus verdaderas ideas, que escondía para evitar problemas con la Iglesia.}
  {Porque así copiaba el orden en que sus estudiantes le hacían las preguntas en clase.}
  {Porque el método exige examinar en serio los mejores argumentos en contra antes de responder.}
\justificacion{La disputa escolástica obliga a considerar las objeciones y a responderlas. A
  ignora que luego viene un «Respondo» claro; B y C inventan motivos que el esquema no permite.}
\especificacion{Interpretativa}{Explica la lógica de un método de argumentación.}
\end{pregunta}

\begin{pregunta}[A]
Una estudiante quiere escribir un ensayo sobre «¿Debe usarse uniforme en el colegio?» con el
método de Tomás. Las partes de su borrador están desordenadas:
\secuencia{«Respondo»: mi postura y mis razones.}{«Parece que no»: los mejores argumentos contra el uniforme.}
  {Respuesta a cada uno de los argumentos contra el uniforme.}{«Pero en contra»: lo que dice el manual de convivencia.}
¿En qué orden debe presentarlas?
\opciones{2-4-1-3}{1-2-3-4}{4-2-1-3}{2-1-4-3}
\justificacion{Objeciones (2), autoridad en contra (4), respuesta propia (1), respuestas a las
  objeciones (3). B pone primero la conclusión; C y D alteran el lugar de la autoridad.}
\especificacion{Propositiva}{Aplica un método de argumentación a un problema nuevo.}
\end{pregunta}

\begin{pregunta}[A]
\afirmacion{en la escolástica, la filosofía de Aristóteles se usó para pensar con rigor temas de
la fe cristiana}{las obras de Aristóteles, recuperadas gracias a traducciones del árabe y del
griego, ofrecían herramientas lógicas para argumentar}
\afirmacionrazon
\justificacion{Ambas son verdaderas y la razón explica la afirmación: la lógica aristotélica, que
  volvió a Europa con las traducciones, dio a la escolástica su método de argumentación.}
\especificacion{Argumentativa}{Explica la relación entre la recuperación de Aristóteles y la escolástica.}
\end{pregunta}

% ================================================================ Bloque 4: ¿Edad Oscura?
\begin{contexto}[¿Fue «oscura» la Edad Media?]
\grupo{Responde las preguntas 9 y 10 con base en el texto y la gráfica.}
El poeta italiano Petrarca (s.\ XIV) llamó «tinieblas» a los siglos que siguieron a la caída de
Roma, y en los siglos siguientes se habló de una «Edad Oscura» o de «oscurantismo». Es cierto que
hubo guerras, epidemias, censura y que la mayoría de la población no sabía leer. Pero en esos
siglos los monasterios copiaron a mano los textos antiguos, se construyeron catedrales y nacieron
las universidades.
\begin{center}
\begin{tikzpicture}
\begin{axis}[eval,ybar,height=3.9cm,width=0.92\linewidth,ymin=0,ymax=42,bar width=16pt,
  symbolic x coords={XI--XII,XIII,XIV,XV},xtick=data,nodes near coords,
  ylabel={Universidades},xlabel={Siglo de fundación},enlarge x limits=0.18]
  \addplot coordinates {(XI--XII,5) (XIII,18) (XIV,22) (XV,35)};
\end{axis}
\end{tikzpicture}
\par{\fontsize{6.5}{7.5}\selectfont\color{black!65}Universidades fundadas en Europa por siglo (cifras aproximadas).\par}
\end{center}
\end{contexto}

\begin{pregunta}[D]
¿Cuál es la conclusión \textbf{más rigurosa} sobre llamar «oscura» a toda la Edad Media?
\opcioneslargas
  {Que fue una época en la que no existió ningún tipo de conocimiento ni de estudio.}
  {Que fue una época perfecta, sin censura, sin guerras y con educación para toda la población.}
  {Que el nombre «Edad Oscura» lo inventaron los propios medievales para describir su tiempo.}
  {Que es una etiqueta discutible: hubo crisis y censura, pero también copia de textos y universidades.}
\justificacion{D reconoce los problemas y también los aportes que el texto y la gráfica muestran.
  A y B son extremos opuestos; C es falso: la etiqueta la pusieron autores posteriores, como Petrarca.}
\especificacion{Argumentativa}{Evalúa una etiqueta histórica con evidencias a favor y en contra.}
\end{pregunta}

\begin{pregunta}[B]
¿Qué muestra la gráfica?
\opcioneslargas
  {Que no se fundó ninguna universidad en Europa antes de que comenzara el Renacimiento.}
  {Que la fundación de universidades creció desde el siglo XIII, antes del Renacimiento.}
  {Que el número de universidades fundadas disminuyó con el paso de los siglos.}
  {Que en los siglos XI y XII se fundaron más universidades que en cualquier otro siglo.}
\justificacion{Se pasa de 5 a 18, 22 y 35 fundaciones. A contradice los datos (hay universidades
  desde el siglo XI); C invierte la tendencia; D lee mal la barra más baja.}
\especificacion{Interpretativa}{Lee una gráfica de barras y extrae una tendencia.}
\end{pregunta}

% ================================================================ Bloque 5: Renacimiento y Erasmo
\begin{contexto}[Renacimiento, imprenta y humanismo]
\grupo{Responde las preguntas 11 y 12 con base en los textos y las imágenes.}
\retrato[0.36]{img/filosofia10/imprenta.jpg}{Taller de imprenta, grabado de Jost Amman (1568).}%
{\small Hacia 1450, Gutenberg perfeccionó la imprenta de tipos móviles. Un libro que antes se
copiaba a mano durante meses podía producirse en cientos de ejemplares. Los \textbf{humanistas}
quisieron volver a leer directamente a los autores griegos y latinos (su lema: «\emph{ad fontes}»,
«a las fuentes») y pusieron en el centro la dignidad y las capacidades del ser humano, sin dejar de
ser, en su mayoría, cristianos.}
\retrato[0.27]{img/filosofia10/erasmo.jpg}{Erasmo de Róterdam (1466--1536), según Hans Holbein.}%
{\small Erasmo fue el humanista más leído de su tiempo. En el \emph{Elogio de la locura} (1511),
dedicado a su amigo Tomás Moro, la Locura habla en primera persona y se burla de los teólogos que
discuten sin fin sobre sutilezas, de los poderosos vanidosos y de la superstición. El título latino,
\emph{Moriae encomium}, juega con «\emph{moría}» (locura, en griego) y el apellido de Moro.}
\end{contexto}

\begin{pregunta}[C]
¿Por qué la imprenta fue tan importante para el humanismo?
\opcioneslargas
  {Porque obligó a dejar de leer los textos antiguos y a escribir solo libros nuevos.}
  {Porque al principio solo se usó para imprimir Biblias, y los humanistas querían leerlas.}
  {Porque permitió multiplicar los libros, y así las ideas de autores como Erasmo circularon por Europa.}
  {Porque fue inventada por Erasmo para poder publicar el \emph{Elogio de la locura}.}
\justificacion{La imprenta abarató y multiplicó los libros: las ideas humanistas se difundieron
  con rapidez. A contradice el lema \emph{ad fontes}; B limita la imprenta a un solo uso; D atribuye
  el invento a quien no fue.}
\especificacion{Argumentativa}{Relaciona un cambio técnico con la difusión de las ideas.}
\end{pregunta}

\begin{pregunta}[A]
Según el texto, ¿qué caracteriza al humanismo frente a la escolástica de las universidades?
\opcioneslargas
  {Volver a leer a los clásicos en sus fuentes y poner en el centro la dignidad del ser humano.}
  {Negar la existencia de Dios y abandonar por completo cualquier forma de religión.}
  {Retirarse a los monasterios y alejarse del mundo para dedicarse solo a la oración.}
  {Dejar de estudiar el latín y el griego para escribir únicamente en las lenguas locales.}
\justificacion{\emph{Ad fontes} y la dignidad humana definen al humanismo. B es un error frecuente
  (casi todos los humanistas eran cristianos); C describe la vida monástica; D contradice el
  regreso a las fuentes griegas y latinas.}
\especificacion{Interpretativa}{Caracteriza un movimiento intelectual frente a otro.}
\end{pregunta}

% ================================================================ Bloque 6: la Utopía
\begin{contexto}[Tomás Moro y la \emph{Utopía} (1516)]
\grupo{Responde las preguntas 13 a 16 con base en los textos y las imágenes.}
\retrato[0.27]{img/filosofia10/moro.jpg}{Tomás Moro (1478--1535), grabado según Hans Holbein.}%
{\small Tomás Moro fue abogado, humanista y canciller del rey Enrique VIII de Inglaterra. Fue
condenado a muerte en 1535 por negarse a reconocer al rey como jefe de la Iglesia inglesa. En la
\emph{Utopía}, el viajero Rafael Hitlodeo critica primero la Inglaterra de su tiempo (libro I) y
luego describe una isla lejana, Utopía (libro II). El nombre viene del griego \emph{ou-tópos}
(«no-lugar») y suena igual que \emph{eu-tópos} («buen lugar»).}
\figura[0.98]{img/filosofia10/utopia.jpg}{Primera edición de la \emph{Utopía} (Lovaina, 1516):
mapa de la isla y el alfabeto inventado de los utopianos, en un ejemplar coloreado a mano.}
\fragmento{Vuestras ovejas, que solían ser tan mansas y comer tan poco, se han vuelto tan voraces y
salvajes que devoran a los propios hombres; devastan y despueblan campos, casas y aldeas.}{Moro,
\emph{Utopía}, libro I, adaptado}
\par\vspace{2pt}{\fontsize{7.5}{9}\selectfont\renewcommand{\arraystretch}{1.2}%
\begin{tabularx}{\linewidth}{|>{\raggedright\arraybackslash}X|>{\raggedright\arraybackslash}X|}\hline
\rowcolor{acento!15}\textbf{Inglaterra hacia 1516} & \textbf{La isla de Utopía}\\\hline
Los nobles cercan los campos para criar ovejas, porque la lana deja más ganancia; los campesinos
quedan sin tierra. & No hay propiedad privada: los bienes son comunes y cada familia toma lo que
necesita.\\\hline
Muchos trabajan de sol a sol y otros viven sin trabajar. & Todos trabajan, unas seis horas al día;
el resto del tiempo es para aprender.\\\hline
El oro y las joyas son señal de poder y causa de envidia. & El oro se usa para las bacinillas y las
cadenas de los esclavos.\\\hline
\end{tabularx}}
\end{contexto}

\begin{pregunta}[C]
¿Qué significado tiene el nombre «Utopía», según el texto?
\opcioneslargas
  {Es el nombre de una isla real del Atlántico que Moro conoció durante sus viajes.}
  {En latín significa «ciudad de Dios», en homenaje a la obra de san Agustín.}
  {Juega con «no-lugar» y «buen lugar»: una sociedad ideal que no existe en ninguna parte.}
  {Era el nombre de la casa de campo de Moro, donde escribió el libro con Erasmo.}
\justificacion{El juego entre \emph{ou-tópos} y \emph{eu-tópos} expresa una sociedad buena que no
  existe. A y D inventan referencias reales; B confunde el nombre con la obra de Agustín.}
\especificacion{Interpretativa}{Explica el sentido de un término a partir de su origen.}
\end{pregunta}

\begin{pregunta}[C]
¿Por qué en Utopía el oro se usa para hacer bacinillas y cadenas de esclavos?
\opcioneslargas
  {Porque en la isla el oro era tan abundante que no servía para ninguna otra cosa útil.}
  {Porque el oro es un metal más resistente y duradero que el hierro o el cobre.}
  {Para que nadie lo desee ni lo admire, y así no haya codicia ni desigualdad entre las personas.}
  {Para poder venderlo después a otros países a un precio más alto y enriquecer a la isla.}
\justificacion{Asociar el oro a objetos despreciables quita su prestigio y con él la codicia, que
  en Inglaterra era causa de envidia. A y B dan razones materiales que el texto no menciona; D
  contradice que en Utopía no haya afán de riqueza.}
\especificacion{Argumentativa}{Infiere la intención de una medida de la sociedad ideal.}
\end{pregunta}

\begin{pregunta}[B]
¿Qué critica Moro con la imagen de las ovejas que «devoran a los hombres»?
\opcioneslargas
  {Que en Inglaterra las ovejas se habían vuelto animales peligrosos que atacaban a los campesinos.}
  {Que los nobles convertían los campos en pastos para la lana y dejaban a los campesinos en la miseria.}
  {Que los ingleses comían demasiada carne de oveja y por eso había hambre en los pueblos.}
  {Que una plaga había llegado a Inglaterra en los rebaños de ovejas traídos de otros países.}
\justificacion{Es una metáfora del cercamiento de tierras (ver la tabla): la ganancia de la lana
  expulsa a los campesinos. A es la lectura literal; C y D proponen causas que el texto no da.}
\especificacion{Interpretativa}{Interpreta una metáfora como crítica social.}
\end{pregunta}

\begin{pregunta}[B]
Un estudiante dice: «La \emph{Utopía} no sirve para nada, porque esa isla no existe». ¿Cuál es la
mejor respuesta?
\opcioneslargas
  {Tiene razón: un libro solo sirve si describe lugares reales que se puedan visitar.}
  {Una utopía no describe un lugar real: muestra por contraste los defectos de la sociedad y permite imaginar cómo mejorarla.}
  {La isla sí existe: es la misma Inglaterra, que Moro describió con otro nombre para no ser castigado.}
  {Moro quería engañar a sus lectores para que viajaran a buscar una isla que él sabía que no existía.}
\justificacion{La comparación de la tabla muestra la función crítica de la utopía. A reduce el valor
  de un texto a lo real; C confunde la crítica con una descripción en clave; D atribuye una
  intención absurda al autor.}
\especificacion{Argumentativa}{Evalúa la función de un texto utópico.}
\end{pregunta}

% ================================================================ Abiertas
\begin{pregunta}[abierta]
Completa la tabla con \textbf{tres} de estas figuras: Agustín, Tomás de Aquino, Erasmo, Moro.
\tablarellenar{3}{Figura}{¿A quién o a qué se le debe? (autor griego, contexto histórico)}{Su aporte principal}
\modelo{Agustín: Platón y los neoplatónicos; crisis del Imperio romano; verdad interior, fe y razón.
Tomás: Aristóteles recuperado por traducciones; universidades; \emph{Suma}, método escolástico.
Erasmo: humanismo, imprenta, \emph{ad fontes}; crítica en el \emph{Elogio de la locura}. Moro:
humanismo, Inglaterra de los cercamientos; la \emph{Utopía}. Criterios: 1 punto por casilla bien
completada (máximo 6).}
\end{pregunta}

\begin{pregunta}[abierta]
Siguiendo el ejemplo de Moro, diseña una «utopía» para tu colegio o tu país: escribe \textbf{tres}
reglas de tu sociedad ideal y, para cada una, el problema actual que estás criticando.
\abierta[6]
\modelo{Se valora: tres reglas claras (3) y, para cada una, un problema real que critica, como hace
la tabla de Moro (3). Se valora más si alguna regla, como el oro en Utopía, cambia el valor que se le
da a algo.}
\end{pregunta}

\begin{pregunta}[abierta]
¿Es justo llamar «Edad Oscura» a la Edad Media? Toma una postura y apóyala con \textbf{dos}
evidencias de esta evaluación.
\abierta[6]
\modelo{Respuesta abierta. Criterios: postura clara (1); dos evidencias tomadas de la evaluación,
como universidades, copia de textos, escolástica, traducciones, o censura y analfabetismo (4);
reconocimiento de lo que diría la postura contraria (1).}
\end{pregunta}

\begin{pregunta}[abierta]
Si Erasmo escribiera hoy un \emph{Elogio de la locura}, ¿de qué «locuras» de nuestra sociedad se
burlaría? Escribe un párrafo breve en el que hable la Locura en primera persona.
\abierta[6]
\modelo{Se valora: uso de la voz de la Locura en primera persona, como en Erasmo (2); crítica de
al menos dos comportamientos actuales, como la vanidad en redes, la superstición o las discusiones
vacías (2); ironía bien lograda (2).}
\end{pregunta}

\end{preguntas}

\autoevaluacion
\creditos{Textos: Agustín, \emph{Confesiones}, \emph{La verdadera religión} y \emph{Sermón} 43; Tomás
de Aquino, \emph{Suma teológica}; Erasmo, \emph{Elogio de la locura} (1511); Moro, \emph{Utopía}
(1516); traducciones adaptadas. Imágenes de dominio público: \emph{La escuela de Atenas} (Rafael),
Wikimedia Commons; Agustín, Tomás de Aquino, Erasmo (según Holbein), Moro (J. Houbraken según
Holbein) e imprenta (J. Amman), Wellcome Collection. Mapa y alfabeto de Utopía (1516), vía Luminarium (luminarium.org).}
\tablaespecificaciones
\imprimirclave
\end{document}
